\documentclass[10pt,a4paper]{amsart}
\usepackage{amsmath,amssymb,amsthm,graphicx}
% Retained minimal preamble for the Stacks Project chapter more-morphisms.tex (More on Morphisms).
% Verbatim excerpt of the chapter preamble: Xy-pic setup (native preamble lines 14-18) and the native
% theorem declarations with the mathematical macros (native preamble lines 154-217). Copied mechanically;
% nothing added, removed or redefined. The wrapper class and the amsmath/amssymb/amsthm/graphicx base
% packages are supplied by the assembler, so the class/externaldocument part of the native preamble is omitted.
% For commutative diagrams we use Xy-pic
\usepackage[all]{xy}

% We use 2cell for 2-commutative diagrams.
\xyoption{2cell}
% Theorem environments.
%
\theoremstyle{plain}
\newtheorem{theorem}[subsection]{Theorem}
\newtheorem{proposition}[subsection]{Proposition}
\newtheorem{lemma}[subsection]{Lemma}

\theoremstyle{definition}
\newtheorem{definition}[subsection]{Definition}
\newtheorem{example}[subsection]{Example}
\newtheorem{exercise}[subsection]{Exercise}
\newtheorem{situation}[subsection]{Situation}

\theoremstyle{remark}
\newtheorem{remark}[subsection]{Remark}
\newtheorem{remarks}[subsection]{Remarks}

\numberwithin{equation}{subsection}

% Macros
%
\def\lim{\mathop{\mathrm{lim}}\nolimits}
\def\colim{\mathop{\mathrm{colim}}\nolimits}
\def\Spec{\mathop{\mathrm{Spec}}}
\def\Hom{\mathop{\mathrm{Hom}}\nolimits}
\def\Ext{\mathop{\mathrm{Ext}}\nolimits}
\def\SheafHom{\mathop{\mathcal{H}\!\mathit{om}}\nolimits}
\def\SheafExt{\mathop{\mathcal{E}\!\mathit{xt}}\nolimits}
\def\Sch{\mathit{Sch}}
\def\Mor{\mathop{\mathrm{Mor}}\nolimits}
\def\Ob{\mathop{\mathrm{Ob}}\nolimits}
\def\Sh{\mathop{\mathit{Sh}}\nolimits}
\def\NL{\mathop{N\!L}\nolimits}
\def\CH{\mathop{\mathrm{CH}}\nolimits}
\def\proetale{{pro\text{-}\acute{e}tale}}
\def\etale{{\acute{e}tale}}
\def\QCoh{\mathit{QCoh}}
\def\Ker{\mathop{\mathrm{Ker}}}
\def\Im{\mathop{\mathrm{Im}}}
\def\Coker{\mathop{\mathrm{Coker}}}
\def\Coim{\mathop{\mathrm{Coim}}}

% Boxtimes
%
\DeclareMathSymbol{\boxtimes}{\mathbin}{AMSa}{"02}

%
% Macros for moduli stacks/spaces
%
\def\QCohstack{\mathcal{QC}\!\mathit{oh}}
\def\Cohstack{\mathcal{C}\!\mathit{oh}}
\def\Spacesstack{\mathcal{S}\!\mathit{paces}}
\def\Quotfunctor{\mathrm{Quot}}
\def\Hilbfunctor{\mathrm{Hilb}}
\def\Curvesstack{\mathcal{C}\!\mathit{urves}}
\def\Polarizedstack{\mathcal{P}\!\mathit{olarized}}
\def\Complexesstack{\mathcal{C}\!\mathit{omplexes}}
% \Pic is the operator that assigns to X its picard group, usage \Pic(X)
% \Picardstack_{X/B} denotes the Picard stack of X over B
% \Picardfunctor_{X/B} denotes the Picard functor of X over B
\def\Pic{\mathop{\mathrm{Pic}}\nolimits}
\def\Picardstack{\mathcal{P}\!\mathit{ic}}
\def\Picardfunctor{\mathrm{Pic}}
\def\Deformationcategory{\mathcal{D}\!\mathit{ef}}
\usepackage{cleveref}
\crefname{originallemma}{original lemma}{original lemmas}
\Crefname{originallemma}{Original lemma}{Original lemmas}
\makeatletter
\newlabel{lemma-composition-perfect}{{0689}{source}{}{sourcepointer-1}{}}
\newlabel{lemma-composition-perfect@cref}{{[originallemma][0][]0689}{[1][0][]source}{}{}{}}
\makeatother
\crefname{originallemma}{original lemma}{original lemmas}
\Crefname{originallemma}{Original lemma}{Original lemmas}
\makeatletter
\newlabel{lemma-flat-finite-presentation-perfect}{{068A}{source}{}{sourcepointer-2}{}}
\newlabel{lemma-flat-finite-presentation-perfect@cref}{{[originallemma][0][]068A}{[1][0][]source}{}{}{}}
\makeatother
\crefname{originallemma}{original lemma}{original lemmas}
\Crefname{originallemma}{Original lemma}{Original lemmas}
\makeatletter
\newlabel{lemma-regular-immersion-perfect}{{068C}{source}{}{sourcepointer-3}{}}
\newlabel{lemma-regular-immersion-perfect@cref}{{[originallemma][0][]068C}{[1][0][]source}{}{}{}}
\makeatother
\title{A local complete intersection morphism is locally of finite presentation, pseudo-coherent and perfect}
\author{The Stacks Project Authors}
\date{stacks-project:more-morphisms.tex:local-snapshot-20261011; GNU FDL 1.2 or later}
\begin{document}
\maketitle
\noindent\textit{Editorial note.} Counted span is the author's own proof environment, source lines 18160-18175, immediately following the statement at 18149-18158.\par\medskip
\noindent\textit{Editorial note.} Same-name twin check: the identical author label in spaces-more-morphisms (statement 11207-11216) has an omitted proof and is therefore not counted; this card is the schemes version with the complete author proof.\par\medskip
\noindent\textit{Editorial note.} Three chapter lemmas are cited by \textbackslash{}ref as navigation pointers; the reduction chain that uses them is inside the counted span.\par\medskip
\noindent\textit{Allowed standard prerequisites (editorial).}\begin{itemize}
\item A Koszul-regular immersion is perfect. chapter lemma lemma-regular-immersion-perfect, cited at proof line 18169. Supplies perfectness of the closed immersion i in the local factorisation of the morphism.
\item A flat morphism of finite presentation is perfect. chapter lemma lemma-flat-finite-presentation-perfect, cited at proof line 18172. Supplies perfectness of the smooth morphism pi in the factorisation.
\item Composition of perfect morphisms is perfect. chapter lemma lemma-composition-perfect, cited at proof line 18174. Combines the two pieces into perfectness of the composed morphism.
\item Perfect ring maps are pseudo-coherent and pseudo-coherent ring maps are of finite presentation. used at proof lines 18161-18165. Derives claims (1) and (2) from claim (3), reducing the whole lemma to perfectness.
\item Local nature of perfectness. used at proof line 18165. Allows the defining local factorisation into a smooth morphism followed by a Koszul-regular immersion to be assumed.
\end{itemize}\medskip
\ifdefined\hypertarget\hypertarget{sourcepointer-1}{}\fi
\noindent\textit{Source reference.} Original lemma 0689: sources/stacks-more-morphisms/pinned-original-source.tex, source lines 17760--17760. Navigation pointer only: the counted proof cites this statement by \textbackslash{}ref as a standing input; the deduction the counted proof has to make is carried out inside the counted span, and no part of that deduction is delegated to this pointer.\par\medskip
\ifdefined\hypertarget\hypertarget{sourcepointer-2}{}\fi
\noindent\textit{Source reference.} Original lemma 068A: sources/stacks-more-morphisms/pinned-original-source.tex, source lines 17783--17783. Navigation pointer only: the counted proof cites this statement by \textbackslash{}ref as a standing input; the deduction the counted proof has to make is carried out inside the counted span, and no part of that deduction is delegated to this pointer.\par\medskip
\ifdefined\hypertarget\hypertarget{sourcepointer-3}{}\fi
\noindent\textit{Source reference.} Original lemma 068C: sources/stacks-more-morphisms/pinned-original-source.tex, source lines 17814--17814. Navigation pointer only: the counted proof cites this statement by \textbackslash{}ref as a standing input; the deduction the counted proof has to make is carried out inside the counted span, and no part of that deduction is delegated to this pointer.\par\medskip
\begin{lemma}
\label{lemma-lci-properties}
Let $f : X \to S$ be a local complete intersection morphism.
Then
\begin{enumerate}
\item $f$ is locally of finite presentation,
\item $f$ is pseudo-coherent, and
\item $f$ is perfect.
\end{enumerate}
\end{lemma}

\begin{proof}
Since a perfect morphism is pseudo-coherent
(because a perfect ring map is pseudo-coherent)
and a pseudo-coherent morphism is locally of finite presentation
(because a pseudo-coherent ring map is of finite presentation)
it suffices to prove the last statement. Being perfect is a local
property, hence we may assume that $f$ factors as $\pi \circ i$ where
$\pi$ is smooth and $i$ is a Koszul-regular immersion.
A Koszul-regular immersion is perfect, see
Lemma \ref{lemma-regular-immersion-perfect}.
A smooth morphism is perfect as it is flat and locally of finite
presentation, see
Lemma \ref{lemma-flat-finite-presentation-perfect}.
Finally a composition of perfect morphisms is perfect, see
Lemma \ref{lemma-composition-perfect}.
\end{proof}

\end{document}
